\documentclass[11pt, a4paper]{article}
\usepackage[a4paper, top=2.5cm, bottom=2.5cm, left=2cm, right=2cm]{geometry}
\usepackage{fontspec}
\usepackage{amsmath, amssymb, amsthm}
\usepackage{listings}
\usepackage{xcolor}
\usepackage{hyperref}
\usepackage{enumitem}
\usepackage{graphicx}
\usepackage[english]{babel}

\babelprovide[import, onchar=ids fonts]{english}
\babelfont{rm}{Noto Sans}

% Define colors for code listings
\definecolor{codegreen}{rgb}{0,0.6,0}
\definecolor{codegray}{rgb}{0.5,0.5,0.5}
\definecolor{codepurple}{rgb}{0.58,0,0.82}
\definecolor{backcolour}{rgb}{0.95,0.95,0.92}

\lstdefinestyle{leanstyle}{
    backgroundcolor=\color{backcolour},
    commentstyle=\color{codegreen},
    keywordstyle=\color{magenta},
    numberstyle=\tiny\color{codegray},
    stringstyle=\color{codepurple},
    basicstyle=\ttfamily\footnotesize,
    breakatwhitespace=false,
    breaklines=true,
    captionpos=b,
    keepspaces=true,
    numbers=left,
    numbersep=5pt,
    showspaces=false,
    showstringspaces=false,
    showtabs=false,
    tabsize=2
}

\lstset{style=leanstyle}

\title{\textbf{The Address is the Map: v3.0}\\
\large The Structural Security Paradox: Why $P=NP$ Doesn't Break RSA}
\author{Based on the Electronic Graph Paper Theory (EGPT) Framework}
\date{\today}

\begin{document}

\maketitle

\begin{abstract}
The resolution of the $P$ vs $NP$ problem via \textbf{Electronic Graph Paper Theory (EGPT)} asserts that the information required to \textit{specify} a unique solution (Address) is identical to the work required to \textit{compute} it (Map). This seemingly implies that all "hard" problems, including integer factorization, are trivial. This paper resolves this apparent contradiction. We demonstrate that cryptographic security is not based on "computational hardness" in the abstract Big-O sense, but on \textbf{Structural Security} derived from lossy compression. In standard number theory, the "Map" (prime factorization) is lossily compressed into the "Address" (composite number). Recovering the Map is equivalent to a backwards walk through Pascal's Triangle, where the search space is physically constrained by binomial coefficients. We introduce \textbf{Recursive Logarithmic Decomposition (RLD)}, a practical Shannon Coding scheme, to show that "breaking" crypto is equivalent to asking whether the map data has already been pre-computed and cached.
\end{abstract}

\section{Introduction: The Paradox of Solvability}

If $P=NP$, as proven in the EGPT framework, then finding a solution should be as fast as verifying it. Why, then, is the world's banking system not currently collapsing?

The answer lies in a fundamental misunderstanding of what an "Address" is in standard mathematics versus a constructive physics.
\begin{itemize}
    \item \textbf{The EGPT Proof:} Proves that if you have the full "Address" (the constructive definition of the state), you implicitly have the "Map" (the algorithm to reach it). $P=NP$ is a conservation law of information.
    \item \textbf{The Cryptographic Reality:} Standard number theory (Cantor space) uses a \textbf{Lossy Compression}. It gives you the "Address" (the composite number $N$) but throws away the "Map" (the prime factors $p, q$).
\end{itemize}

Security, therefore, is not a failure of algorithms ($P \neq NP$). It is a success of information hiding. The "hardness" of factoring is the physical work required to regenerate the lost map information.

\section{Deconstructing the "Address is the Map"}

To understand this, we must revisit the central thesis of EGPT with rigor.

\subsection{The "Free Address" Fallacy vs. Physical Addressing}
In the standard Random Access Memory (RAM) model, we assume we can "jump" to memory address $10^{100}$ in $O(1)$ time. This is the **Free Address Fallacy**.
\begin{itemize}
    \item **Standard View:** The "Address" is a scalar label. The "Search" is the work to find the label.
    \item **Physical View (EGPT):** To *define* a unique location in a configuration space of size $2^K$, you must provide $K$ bits of information. This sequence of $K$ bits *is* the path (the Map) to the location.
\end{itemize}

\textbf{Analogy: The Manhattan Grid}
If I give you the address "24 East, 32 North," I have explicitly given you the map: "Walk 24 blocks East, then 32 blocks North." The work is polynomial in the address length.
However, if I compress this address into a single scalar "Building \#5943" using a secret hash function (lossy compression of the path), you can no longer "walk" there. You must search.

The "search" in $NP$ problems is simply the work required to **uncompress** the lossy address back into its coordinate map.

\section{The Ubiquity of Pascal's Triangle}

Why does this decompression process always seem to involve combinatorics? Why do we see the Binomial Theorem in Newton's calculation of $\pi$, in the "Twiddle Factors" of the Fast Fourier Transform (FFT), and in the structure of prime distribution?

\subsection{The Bernoulli Process Foundation}
In EGPT, every number and event is the result of finite sampling from an infinite, Independent and Identically Distributed (IID) source (a "fair coin").
\begin{itemize}
    \item **The Binary Tree:** A random walk creates a binary tree of possibilities.
    \item **The Binomial Coefficient:** The number of ways to reach a specific state (e.g., $k$ heads in $n$ flips) is given by $\binom{n}{k}$, the entries of Pascal's Triangle.
\end{itemize}

\subsection{Newton, Pi, and the Negative Binomial}
Newton calculated $\pi$ efficiently not by magic, but by using the **Negative Binomial Theorem** $(1-x)^{-n}$. He realized that the geometric properties of the circle could be "counted" by expanding this series. This is a direct application of combinatorial probability to continuous geometry.

\subsection{FFT and Roots of Unity}
The Fast Fourier Transform (FFT) is the algorithm of decomposing a signal into its frequencies. The "Twiddle Factors" (roots of unity) that drive the FFT are intimately linked to Pascal's Triangle. They act as "filters" that select specific sums of binomial coefficients (e.g., summing every $m$-th entry).
\begin{itemize}
    \item **Connection:** The FFT is essentially a "change of basis" from the time domain (path history) to the frequency domain (roots of unity).
    \item **EGPT Insight:** Primes act as the "frequencies" of the number system. Factoring is the inverse FFT of the integers.
\end{itemize}

\section{Practical Implementation: Recursive Logarithmic Decomposition (RLD)}

EGPT is not just theory; it is code. We have implemented a practical Shannon Coding scheme called **Recursive Logarithmic Decomposition (RLD)** (formerly PPF) in the `EGPTMath` Javascript library.

\subsection{The Encoding}
RLD provides a constructive bijection between Natural numbers ($\mathbb{N}$) and the Unit Circle (Phase $[0, 1)$).
\begin{lstlisting}[language=C, caption={RLD Canonical Formula}]
// RLD Phase Formula
// N = 2^N_level + offset
phase_RLD = offset / 2^(N_level + 1)
\end{lstlisting}
This formula maps every integer $N$ to a unique phase angle. It unifies:
\begin{itemize}
    \item **Napier's Logarithms:** Recursive halving of the space.
    \item **Shannon's Information:** Binary encoding of the offset.
    \item **Cooley-Tukey FFT:** The even/odd decomposition of the unit circle.
\end{itemize}

\subsection{The "Phase" as Map}
In RLD, the number is not a scalar; it is a **Vector** of phase information. The "Phase" tells you exactly where the number sits in the Bernoulli binary tree. This encoding is **lossless** and **reversible**.

\section{The Structural Security Paradox}

Here we resolve the conflict between $P=NP$ and Crypto.

\subsection{Lossy Compression vs. Lossless Map}
\begin{itemize}
    \item **RLD/EGPT:** Uses **Lossless** Shannon Coding. The prime structure is preserved or easily recoverable because the "Address" (Phase) retains the path information. $P=NP$ is obvious here.
    \item **Standard Math (Cantor):** Uses **Lossy** Compression. The number $N = p \times q$ is stored as a scalar. The "path information" (that it came from $p$ and $q$) is discarded.
\end{itemize}

\subsection{Walking Backwards Through Pascal's Triangle}
"Breaking" RSA is the process of reversing this lossy compression. Without the map (the primes), you are effectively at a node in Pascal's Triangle, trying to determine which path led you there.
\begin{itemize}
    \item **The Ambiguity:** A composite number $N$ has a "high Pascal coefficient"—there are many potential paths (combinations of factors) that *could* construct it.
    \item **The Constraint:** Only *one* of those paths is the "Prime" path.
    \item **The Search Range:** The difficulty of finding the factors is determined by the **Binomial Coefficient** corresponding to that number's magnitude. This defines the "Search Range."
\end{itemize}

\subsection{Has the Map Been Pre-Computed?}
The "hardness" of factoring reduces to a simple question: **"Has this address been mapped before?"**
\begin{itemize}
    \item **Pre-computation:** We consider FFTs "fast" ($N \log N$) because we pre-compute and cache the transcendental "Twiddle Factors" ($\pi$, $e$, roots of unity). We don't count the cost of computing $\pi$ every time we run an FFT.
    \item **Primes as Twiddles:** The prime numbers are the "Twiddle Factors" of the number system. If we have a list of pre-computed primes (a "Log Book"), factoring is fast (Polynomial).
    \item **The Reality:** We *don't* have a complete list of primes for $N=2^{4096}$. The "computation" required to break RSA is essentially the work required to **build the map** to that location for the first time.
\end{itemize}

\section{Conclusion}

$P=NP$ does not mean "Magic." It means **Accountability**.
\begin{enumerate}
    \item In a fully constructive universe (EGPT), you cannot define a problem without defining the solution path. Therefore, $P=NP$.
    \item Cryptography relies on the fact that we use **Lossy Number Systems** (Standard Math) that discard the solution path.
    \item Security is guaranteed not by "computational hardness" (which is an illusion of bad accounting), but by the **Structural Entropy** required to rebuild the lost map.
\end{enumerate}

\end{document}